\subsection{Real system implementation and validation}
Different tunings of $Q_{aug}$ and $R_{aug}$ were evaluated both in simulation and on the real apparatus.
However, the tuning configurations providing satisfactory simulated performance did not always achieve comparable results
on the physical system, and vice versa.
This discrepancy is likely due to model mismatch, significant parameter uncertainties,
unmodelled nonlinearities and time-varying operating conditions affecting the real magnetic levitation system.

\subsubsection{Time-Domain Validation}
The time-domain logs in \texttt{data/03\_position/LQR/time} contain reference and measured position,
Kalman-filter speed, reference and measured current, and position error.

The following sets of weights were tested during the time-domain validation:
\begin{table}[H]
\centering
\caption{LQR with integral action parameters used for time-domain validation.}
\label{tab:lqri_time_parameters}
\begin{tabular}{cccccc}
\toprule
Test file & $Q_1$ & $Q_2$ & $Q_3$ & $Q_4$ & $R$ \\
\midrule
1 & $5000$ & $5$  & $10^{-9}$ & $3\times10^{-4}$ & $0.05$ \\
2 & $5000$ & $5$  & $10^{-9}$ & $3\times10^{-4}$ & $0.07$ \\
3 & $5000$ & $5$  & $10^{-9}$ & $3\times10^{-4}$ & $0.09$ \\
4 & $5000$ & $25$ & $10^{-9}$ & $3\times10^{-4}$ & $0.07$ \\
5--6 & $5000$ & $25$ & $10^{-9}$ & $5\times10^{-4}$ & $0.07$ \\
\bottomrule
\end{tabular}
\end{table}

The table summarizes the practical retuning process. Increasing the velocity weight and adapting the
integral-state weight reduced oscillations on the real setup.

The representative log shows tracking around $7~\mathrm{mm}$, with visible measurement noise and some
spikes. This confirms the need for the Kalman speed estimate.

\begin{figure}[H]
    \centering
    \includegraphics[width=0.88\textwidth]{../img/LQRI_validation/time/ramp.png}
    \caption{LQI time-domain position tracking on the real setup.}
    \label{fig:lqri_time_position}
\end{figure}

\begin{figure}[H]
    \centering
    \includegraphics[width=0.88\textwidth]{../img/LQRI_validation/time/error.png}
    \caption{Position error during LQI time-domain validation.}
    \label{fig:lqri_time_error}
\end{figure}

\begin{figure}[H]
    \centering
    \includegraphics[width=0.88\textwidth]{../img/LQRI_validation/time/current.png}
    \caption{Current tracking during LQI time-domain validation.}
    \label{fig:lqri_time_current}
\end{figure}

\begin{figure}[H]
    \centering
    \includegraphics[width=0.88\textwidth]{../img/LQRI_validation/time/speed.png}
    \caption{Kalman-filter velocity estimate compared with the derivative-based speed signal.}
    \label{fig:lqri_time_speed}
\end{figure}



\subsubsection{Frequency-Domain Validation}
The frequency-domain validation is performed with sinusoidal references at
\begin{equation}
    \omega = 1,\ 6,\ 15\ \mathrm{rad/s}.
\end{equation}

The measured sinusoidal response is compared with the predicted frequency response in terms of
attenuation and phase shift.

\begin{figure}[H]
    \centering
    \begin{subfigure}[t]{0.48\textwidth}
        \centering
        \includegraphics[width=\textwidth]{../img/LQRI_validation/simulation/LQR_1_rad.png}
        \caption{Simulation, reference $\omega = 1$ rad/s}
        \label{fig:lqr_1_rad_simulation}
    \end{subfigure}
    \hfill
    \begin{subfigure}[t]{0.48\textwidth}
        \centering
        \includegraphics[width=\textwidth]{../img/LQRI_validation/frequency/LQR_1_rad.png}
        \caption{Experimental response, reference $\omega = 1$ rad/s}
        \label{fig:lqr_1_rad}
    \end{subfigure}
    \caption{LQRI validation for a sinusoidal reference with $\omega = 1$ rad/s.}
    \label{fig:lqri_validation_1_rad}
\end{figure}

\begin{figure}[H]
    \centering
    \begin{subfigure}[t]{0.48\textwidth}
        \centering
        \includegraphics[width=\textwidth]{../img/LQRI_validation/simulation/LQR_6_rad.png}
        \caption{Simulation, reference $\omega = 6$ rad/s}
        \label{fig:lqr_6_rad_simulation}
    \end{subfigure}
    \hfill
    \begin{subfigure}[t]{0.48\textwidth}
        \centering
        \includegraphics[width=\textwidth]{../img/LQRI_validation/frequency/LQR_6_rad.png}
        \caption{Experimental response, reference $\omega = 6$ rad/s}
        \label{fig:lqr_6_rad}
    \end{subfigure}
    \caption{LQRI validation for a sinusoidal reference with $\omega = 6$ rad/s.}
    \label{fig:lqri_validation_6_rad}
\end{figure}


